\section{Complex Numbers}

\bx{ %0.7.1
    If $z$ is the complex number $a +ib$, which of the following are synonyms? 
    \eg{4}{
        \task\ modulus of $z$
        \task\ $|z|$
        \task\ $\bar z$
        \task\ $\mathrm{Re} \ z$
        \task\ absolute value of $z$
        \task\ $\mathrm{Im} \ z$
        \task\ complex conjugate of $z$
        \task\ $a$
        \task\ $b$
        \task\ real part of $z$
        \task\ imaginary part of $z$
    }
    \sol{
        \[ \text{modulus of } z \equiv |z| \equiv \text{ absolute value of } z \]
        \[ \bar z \equiv \text{ complex conjugate of } z \]
        \[ \mathrm{Re} \ z \equiv a \equiv \text{ real part of } z \]
        \[ \mathrm{Im} \ z \equiv b \equiv \text{ imaginary part of } z \]
    }
}

\bx{ %0.7.2
    If $z = a + ib$, which of the following are real numbers?
    \eg{4}{
        \task\ $|z|$
        \task\ $\mathrm{Re} \ z$
        \task\ $\sqrt{z \bar z}$
        \task\ $\bar z$
        \task\ $a$
        \task\ $\mathrm{Im} \ z$
        \task\ $b$
        \task\ $z + \bar z$
        \task\ modulus of $z$
        \task\ real part of $z$
        \task\ imaginary part of $z$
        \task\ $z - \bar z$
    }
    \sol{
        The following elements are real numbers
        \[
        |z|, \quad \mathrm{Re} \ z, \quad \sqrt{z \bar z}, \quad a, \quad \mathrm{Im} \ z, \quad b, \quad z + \bar z, \quad \text{ modulus of } z, \quad \text{ real part of } z, \quad \text{imaginary part of } z
        \]
    }
}
\pagebreak
\bx{ %0.7.3
    Find the absolute value and argument of each of the following complex numbers:
    \eah{
        \item $2 + 4i$
        \item $(3 + 4i)^{-1}$
        \item $(1 + i)^5$
        \item $1 + 4 i$
    }
    \sol{
        \ea{
            \item $|2 + 4i| = \sqrt{2^2 + 4^2} = 2\sqrt{5}, \quad \arg (2 + 4i) = \tan^{-1} \left( \frac 42 \right) = \tan^{-1} 2$.
            \item $|(3+4i)^{-1}| = |3+4i|^{-1} = \sqrt{3^2 + 4^2}^{-1} = 1/2, \quad \arg ((3+4i)^{-1}) = -\arg (3+4i) = -\tan^{-1} \left(\frac 43 \right)$
            \item $|(1+i)^5| = |1+i|^5 = \sqrt{2}^5 = 4 \sqrt 2,\quad \arg((1+i)^{5})=5\arg(1+i) = 5 \tan^{-1} 1 = \frac{5\pi}4$
            \item $|1+4i| = \sqrt{1^2+4^2} = \sqrt{17}, \quad \arg (1+4i) = \tan^{-1} 4$
        }
    }
}

\bx{ %0.7.4
    Find the modulus and polar angle of the following complex numbers:
    \eah{
        \item $3 + 2i$
        \item $(1 - i)^4$
        \item $2 + i$
        \item $\sqrt[7]{3 + 4i}$
    }
    \sol{
        \ea{
            \item $|3+2i| = \sqrt{3^2 + 2^2} = \sqrt{13}, \quad \arg(3+2i) = \tan^{-1} \left( \frac23 \right)$.
            \item $|(1-i)^4| = |1-i|^4 = 4, \quad \arg ((1-i)^4) = 4 \arg(1-i) = 4\tan^{-1} (-1) = -\pi \equiv \pi \pmod{2\pi}$.
            \item $|2+i| = \sqrt{2^2 +1^2} = \sqrt 5, \quad \arg (2+i) = \tan^{-1} \left(\frac12 \right)$.
            \item $|\sqrt[7]{3+4i}| = |3+4i|^{1/7} = \sqrt{3^2+4^2}^{1/7} = \sqrt[7]{5}, \quad \arg (\sqrt[7](3+4i)) = \frac17 \tan^{-1}\left( \frac 43\right)$.
        }
    }
}

\bx{ %0.7.5
    Verify the nine rules for addition and multiplication of complex numbers. Statements 5 and 9 are the only ones that are not immediate.
    \sol{
        To recall, the nine rules, for $z_1,z_2,z_3 \in \mathbb C$ are:
        \[
        \begin{cases*}
            1. \ (z_1 + z_2) + z_3 = z_1 + (z_2 + z_3) & Addition is associative. \\
            2. \ z_1 + z_2 = z_2 + z_1 & Addition is commutative. \\
            3. \ z + 0 = z & 0 is an additive identity. \\
            4. \ (a + ib) + (-a - ib) = 0 & Existence of an additive inverse. \\
            5. \ (z_1z_2)z_3 = z_1 (z_2z_3) & Multiplication is associative. \\
            6. \ z_1z_2 = z_2z_1 & Multiplication is commutative. \\
            7. \ 1z = z & 1 is a multiplicative identity. \\
            8. \ (a + ib)\left( \frac{a}{a^2+b^2} - i \frac{b}{a^2+b^2} \right) = 1 & Existence of a multiplicative inverse. \\
            9. \ z_1 (z_2+z_3) = z_1z_2 + z_1z_3 & Multiplication is distributive over addition.
        \end{cases*}
        \]
        Rules 1 to 4 are trivial since the addition of complex numbers is equivalent to the addition under $\R^2$.

        We will show the multiplicative properties. Rule 6 is true since $\R$ also satisfies the multiplicative commutative property
        \[
        (a_1 + ib_1)(a_2 + ib_2) = a_1a_2 - b_1b_2 + i(a_1b_2 + a_2b_1) = a_2a_1 - b_2b_1 + i (b_2a_1 + b_1a_2) = (a_2 + ib_2)(a_1 + ib_1).
        \]
        Rule 7 is trivially true as
        \[
        1(a + ib) = 1a + i(1b) = a + ib.
        \]

        Rule 8 is also true by simply expanding the expression. This is trivial and hence left as an exercise for the reader.

        We have left Rules 5 and 9 to the last since they are more complicated. For Rule 5, we expand the expressions as follows:
        \begin{align*}
        (z_1z_2)z_3
        &=\Bigl[(a_1a_2-b_1b_2)+i(a_1b_2+a_2b_1)\Bigr](a_3+ib_3) \\
        &=\Bigl[(a_1a_2-b_1b_2)a_3-(a_1b_2+a_2b_1)b_3\Bigr]
        +i\Bigl[(a_1a_2-b_1b_2)b_3+(a_1b_2+a_2b_1)a_3\Bigr] \\
        &=(a_1a_2a_3-a_3b_1b_2-a_1b_2b_3-a_2b_1b_3)
        +i(a_1a_2b_3-b_1b_2b_3+a_1a_3b_2+a_2a_3b_1), \\[1ex]
        z_1(z_2z_3)
        &=(a_1+ib_1)\Bigl[(a_2a_3-b_2b_3)+i(a_2b_3+a_3b_2)\Bigr] \\
        &=\Bigl[a_1(a_2a_3-b_2b_3)-b_1(a_2b_3+a_3b_2)\Bigr]
        +i\Bigl[a_1(a_2b_3+a_3b_2)+b_1(a_2a_3-b_2b_3)\Bigr] \\
        &=(a_1a_2a_3-a_1b_2b_3-a_2b_1b_3-a_3b_1b_2)
        +i(a_1a_2b_3+a_1a_3b_2+a_2a_3b_1-b_1b_2b_3).
        \end{align*}
        which is equivalent since addition over $\R$ is also associative.

        Lastly, Rule 9 follows a similar approach:
        \begin{align*}
        z_1(z_2+z_3)
        &=(a_1+ib_1)\bigl[(a_2+a_3)+i(b_2+b_3)\bigr] \\
        &=\Bigl[a_1(a_2+a_3)-b_1(b_2+b_3)\Bigr]
        +i\Bigl[a_1(b_2+b_3)+b_1(a_2+a_3)\Bigr] \\
        &=(a_1a_2+a_1a_3-b_1b_2-b_1b_3)
        +i(a_1b_2+a_1b_3+a_2b_1+a_3b_1), \\
        z_1z_2+z_1z_3
        &=\Bigl[(a_1a_2-b_1b_2)+i(a_1b_2+a_2b_1)\Bigr]
        +\Bigl[(a_1a_3-b_1b_3)+i(a_1b_3+a_3b_1)\Bigr] \\
        &=(a_1a_2+a_1a_3-b_1b_2-b_1b_3)
        +i(a_1b_2+a_1b_3+a_2b_1+a_3b_1).
        \end{align*}
    }
}
\bx{ %0.7.6
    \eah*{
        \item Solve $x^2 + ix = -1$.
        \item Solve $x^4 + x^2 + 1 = 0$
    }
    \sol{
        \item Moving terms to one side gives us $x^2 + ix + 1 = 0$. By quadratic formula, 
        \[
        x = \frac{-i \pm \sqrt{i^2 - 4(1)(1)}}{2} = \frac{-i \pm \sqrt{5}i}{2}.
        \]
        \item Using quadratic formula,
        \begin{align*}
        x^2
        &= \frac{-1 \pm \sqrt{1-4(1)(1)}}{2} \\
        &= \frac{-1 \pm \sqrt{3}\,i}{2} \\
        &= \cos \frac{2\pi}{3} \pm i\sin \frac{2\pi}{3} \\
        &= e^{\pm 2\pi i/3},
        \end{align*}
        so
        \[
        x=e^{i\theta},
        \qquad
        \theta=\pm\frac{\pi}{3},\ \pm\frac{2\pi}{3}.
        \]
    }
}
\pagebreak
\bx{ %0.7.7
    Describe the set of all complex number $z = x + iy$ such that 
    \eah{
        \item $|z - u| + |z - v| = c$ for $u,v \in \mathbb C$ and $c \in \R$
        \item $|z| < 1 - \mathrm{Re} \ z$
    }
    \sol{\ea{
        \item By properties of conic sections, this set represents an ellipse with the foci of $u$ and $v$ on the Argand Diagram. However, we will need to prove this rigorously with the correct conditions. Let $S = \{ z \in C: |z-u| + |z-v| = c \text{ for } u , v \in \C \text{ and } c \in R \}$.
        
        By Triangle Inequality, 
        \[
        |z-u| + |z-v| \ge |u-v|
        \]
        Hence we have to consider the 3 cases, $c < |u - v|$, $c=|u-v|$ and $c > |u-v|$.

        \textbf{Case 1:} Suppose $c < |u-v|$. By triangle Inequality, we get $|z-u|+|z-v| > c$, which suggest there does not exists any $z \in \C$ that satisfy the set. Hence, $S = \emptyset$.

        \textbf{Case 2:} Suppose $c = |u-v|$. The equality is achieved in the Triangle Inequality, hence, $z$, $u$ and $v$ must be collinear with the condition that $z-u$ and $v -z$ must be positive scalar
        multiples of each other. Hence, $S$ is a line on the Argand Diagram connecting $u$ and $v$. To express this formally, we get 
        \[
        S = \{ z = u + t(v-u): t \in [0,1] \}
        \]

        \textbf{Case 3:} Suppose $c > |u-v|$. This case is the most difficult to prove. Before we begin, let us recall the geometric representation of addition and multiplication of complex numbers. 
        Substituting $z$ with $z + c$ caused the point in the Argand diagram to move $\mathrm{Im} \ c$ units downwards and $\mathrm{Re} \ c$ units leftwards. Substituting $z$ with $z e^{i \theta}$ causes the point to rotate $\theta$ radians clockwise.

        Let $\sigma_1$ be a transformation substituting $z$ with $z + w$, $(z \rightarrow z + w)$, where $w = (u+v)/2$, this results in a new set 
        \begin{align*}
        S_1 = \sigma_1 S 
        &= \{ z \in \C: \left|z - \frac v2 + \frac u2\right| + \left|z -\frac u2 + \frac v2\right| = c \text{ for } u,v\in\C, c \in \R \} \\
        &= \{ z \in \C: \left|z - u_1\right| + \left|z -v_1\right| = c \text{ for } u,v\in\C, c \in \R \}
        \end{align*}

        Without loss of generality, suppose $\mathrm{Re} \ u_1 > \mathrm{Re} \ v_1$. Let $\theta = \arg (u_1 -v_1)$ and $\sigma_2$ be the transformation substituting $z$ with $z e^{i\theta}$, 
        ie,$(z \rightarrow ze^{i\theta})$. This gives us a new set 
        \begin{align*}
        S_2 = \sigma_2 S_1 &= \{ z \in \C: \left|ze^{i\theta} - u_1\right| + \left|ze^{i\theta} -v_1\right| = c \text{ for } u,v\in\C, c \in \R \} \\
        & \{ z \in \C: |e^{i\theta}|\left|z - u_1 e^{-i\theta}\right| + |e^{i\theta}| \left|z -v_1 e^{-i\theta}\right| = c \text{ for } u,v\in\C, c \in \R \} \\
        & \{ z \in \C: \left|z - k \right| + \left|z + k\right| = c \text{ for } u,v\in\C, c \in \R \}
        \end{align*}
        Since both points now lie on real axis and are equidistant to the origin, we can substitute them for $k$ and $-k$.

        Now we have to show that this set is an ellipse.
        \begin{align*}
        |z-k|+|z+k|&=c\\
        \sqrt{(x+k)^2+y^2}+\sqrt{(x-k)^2+y^2}&=c\\
        \sqrt{(x+k)^2+y^2}
        &=c-\sqrt{(x-k)^2+y^2}\\
        (x+k)^2+y^2
        &=c^2-2c\sqrt{(x-k)^2+y^2}+(x-k)^2+y^2\\
        4kx
        &=c^2-2c\sqrt{(x-k)^2+y^2}\\
        2c\sqrt{(x-k)^2+y^2}
        &=c^2-4kx\\
        4c^2\bigl((x-k)^2+y^2\bigr)
        &=(c^2-4kx)^2\\
        4c^2(x^2-2kx+k^2+y^2)
        &=c^4-8c^2kx+16k^2x^2\\
        4c^2x^2+4c^2y^2+4c^2k^2
        &=c^4+16k^2x^2\\
        4(c^2-4k^2)x^2+4c^2y^2
        &=c^2(c^2-4k^2).
        \end{align*}
        

        \item Let $z = a + ib$, then
        \begin{align*}
        |z| &< 1 - \mathrm{Re}\,z \\
        \sqrt{a^2+b^2} &< 1-a \implies a<1, \\
        a^2+b^2 &< (1-a)^2 \\
        b^2 &< 1-2a \\
        a &< \frac{1-b^2}{2}.
        \end{align*}

        This results in the region to the left of the parabola
        \[
        a=\frac{1-b^2}{2},
        \]
        which has focus $(0,0)$ and directrix $a=1$, as shown below.

        \begin{figure}[H]
        \centering
        \begin{tikzpicture}
        \begin{axis}[
            axis lines=center,
            xlabel={$\mathrm{Re}\,z$},
            ylabel={$\mathrm{Im}\,z$},
            xmin=-5, xmax=5,
            ymin=-5, ymax=5,
        ]

        \addplot[
            name path=upper,
            domain=-5:0.5,
            samples=200,
            draw=none
        ]
        {sqrt(1-2*x)};

        \addplot[
            name path=lower,
            domain=-5:0.5,
            samples=200,
            draw=none
        ]
        {-sqrt(1-2*x)};

        \addplot[
            fill=blue!20,
            opacity=0.6,
            draw=none
        ]
        fill between[of=upper and lower];

        \addplot[domain=-5:0.5, samples=200, blue, thick]
        {sqrt(1-2*x)};

        \addplot[domain=-5:0.5, samples=200, blue, thick]
        {-sqrt(1-2*x)};

        \addplot[red, thick] coordinates {(1,-5) (1,5)};

        \node[blue] at (axis cs:-2,3.3) {$a<\frac{1-b^2}{2}$};
        \node[right,red] at (axis cs:1,1) {$a=1$};
        \addplot[only marks, mark=*, red] coordinates {(0,0)};
        \node[above left,red] at (axis cs:0,0) {$(0,0)$};

        \end{axis}
        \end{tikzpicture}
        \caption{Region satisfying $\{z\in\mathbb{C}:|z|<1-\mathrm{Re}\,z\}$.}
        \end{figure}
    }
    }
}

\bx{ %0.7.8
    \eah*{
        \item Solve $x^2 + 2ix - 1 = 0$.
        \item Solve $x^3 - x^2 - x = 2$.
    }
}

\bx{ %0.7.9
    Solve the following equation where $x,y$ are complex numbers:
    \eah{
        \item $x^2 + ix + 2 = 0$
        \item $x^4 + x^2 + 2 = 0$
        \item \(
            \begin{cases}
                ix - (2+i)y = 3 \\
                (1+i)x + iy = 4
            \end{cases}
        \)
    }
}

\bx{ %0.7.10
    Describe the set of all complex numbers $z =  x + iy$ such that 
    \eah{
        \item $\bar z = -z$
        \item $|z - a| = |z - b|$ for given $a,b \in \mathbb C$
        \item $\bar z = z^{-1}$
    }
}

\bx{ % 0.7.11
    \ea{
        \item Describe the loci in $\mathbb C$ given by the following equations: \er{
            \item $\mathrm{Re} \ z = 1$,
            \item $|z| = 3$.
        }
        \item What is the image of each locus under the mapping $z \mapsto z^2$?
        \item What is the image image of each locus under the mapping $z \mapsto z^2$?
    }
}

\bx{ % 0.7.12
    \ea{
        \item Describe the loci in $\mathbb C$ given by the following equations: \er{
            \item $\mathrm{Im} \ z = -|z|$,
            \item $\frac12|z-i| = z$.
        }
        \item What is the image of each locus under the mapping $z \mapsto z^2$?
        \item What is the image image of each locus under the mapping $z \mapsto z^2$?
    }
}

\bx{ % 0.7.13
    \ea{
        \item Find all the cubic roots of $1$.
        \item Find all the $4$th roots of $1$.
        \item Find all the $6$th roots of $1$.
    }
}

\bx{ % 0.7.14
    Show that for any complex numbers $z_1, z_2$, we have $\mathrm{Im}(\bar{z_1} z_2) \le |z_1||z_2|$.
}

\bx{ % 0.7.15
    \ea{
        \item Find all the fifth roots of 1, using a formula that only includes square roots (no trigonometric function).
        \item Use your formula, ruler, and compass to construct a regular pentagon.
    }
}
